% results_design_implications

\section{Design Implications}
\label{sec:design-implications}

We do not recommend adding features generically: feature count is not the
lever satisfaction runs on. Each recommendation ties to a specific reported
result, and every coverage recommendation pairs with its delivery point, since
the two gaps need different remedies.

\subsection{Deliver What Is Already Promised, Before Adding Anything}

The two most-broken promises are also the two most basic: adhan reminders
(broken in 5 of 20 promising applications) and monetisation terms (5 of 16)
(\S\ref{sec:rq6}). Prompting failure is independently the dominant tracker
complaint, at 46\% of the cross-app coding pass (\S\ref{sec:rq1}). An
application that cannot reliably fire a notification cannot support a
prayer-tracking practice built on top of it, whatever else it ships.

\subsection{Treat Prayer-Time Correctness as a Reliability Requirement}

Prayer-time accuracy costs \(-1.112\) stars, the second-costliest of the
complaint aspects we model, and 41.6\% of accuracy complaints arrive inside
four- or five-star reviews (\S\ref{sec:rq2}). A team prioritising defects from
its own rating distribution will systematically under-prioritise this class of
failure. We recommend instrumenting for accuracy defects directly, comparing
computed times against an authoritative reference, rather than inferring
severity from ratings that conceal them.

\subsection{Accommodate Practice, Rather Than Softening Mechanics}
\label{sec:design-accommodate}

Where tracker dissatisfaction carries a normative complaint, it is that the
app mis-models worship, not that its mechanics apply too much pressure
(\S\ref{sec:rq1}). Menstrual handling makes this concrete: it ranks second on
the unmet-need index, is absent from 23 of 26 applications, and in the three
that do ship it, sentiment is no better than in applications without it
(\S\ref{sec:rq4}). This is simultaneously a coverage problem and a quality
problem; closing the coverage gap without fixing reliability would not resolve
it.

\subsection{Surface Less; Do Not Necessarily Contain Less}

The lever available to a design team is what the interface surfaces in a
given moment, not how many capabilities the application contains
(\S\ref{sec:rq5}). A full-featured application with a disciplined interface is
not the problem the bloat framing warns against; a lean application that
surfaces its few features poorly can still earn the complaint.

\subsection{The Unserved Categories}

Four categories fall below a quarter of ecosystem coverage against
demonstrated demand (\S\ref{sec:rq6}). \textbf{Companion hardware} tops the
ranking, absent from 25 of 26 applications, and sentiment toward hardware is
\emph{better} inside the one application shipping a device (\S\ref{sec:rq5}),
so this is demand rather than dissatisfaction with an existing product.
\textbf{Menstrual accommodation} ranks second (\S\ref{sec:design-accommodate}).
\textbf{Travel concession} (\emph{qasr}) has essentially no market presence
(\S\ref{sec:rq3}), and \textbf{calendar integration} rounds out the four.
None of these calls for adding features indiscriminately; each is a
demand-evidenced gap in an ecosystem whose more pressing problem is delivering
what it has already promised.
